% CAISc 2026 submission.
% Download caisc_2026.sty from https://caisc2026.github.io/cfp.html
% then compile: pdflatex paper && bibtex paper && pdflatex paper && pdflatex paper

\documentclass{article}

% Submission (default) — anonymized, line numbers, "Submitted to ..." footer.
% Use [preprint] for a non-anonymous arXiv version; [final] for camera-ready.
\usepackage{caisc_2026}

\usepackage[utf8]{inputenc}
\usepackage[T1]{fontenc}
\usepackage[pdfusetitle,breaklinks=true]{hyperref}
\usepackage{url}
\usepackage{booktabs}
\usepackage{amsfonts}
\usepackage{amsmath}
\usepackage{nicefrac}
\usepackage{microtype}
\usepackage{xcolor}
\usepackage{multirow}
\usepackage{graphicx}

\title{Tool-Error-Induced Next-Action Direction in Agentic LLMs}

% Human authors only per CAISc 2026 / OpenReview authorship policy.
% AI system involvement is documented in the AI Involvement Checklist.
\author{Anonymous Author(s) \\ Anonymous Institution \\ \texttt{[anonymous for review]}}

\begin{document}
\maketitle

% ─────────────────────────────────────────────────────────────────────────────
\begin{abstract}
We report a pre-registered mechanistic interpretability study of
next-action choice after tool failure across three 8B-class language
models spanning agent-training depth: Qwen3-8B, Hermes-3-Llama-3.1-8B,
and Llama-xLAM-2-8b-fc-r.
A linear probe trained on the residual stream at the post-tool-result
token (position~$t^*$) achieves \textbf{90.0\%, 86.0\%, and 90.1\%}
balanced accuracy on the 3-class \{retry-same, modify-args, give-up\}
decision (chance 33\%; balanced shuffle floor 33.8\%, 33.6\%, 33.5\%),
with permutation $p<10^{-4}$ surviving Holm-Bonferroni correction.
The probe outperforms a position-only baseline by 27--33\,pp.
Causal steering at $\alpha{=}2.5$ converts give-up trajectories to
any non-give-up action in \textbf{100\%} (Qwen3-8B; $n{=}60$,
file\_read-dominated---see \S\ref{sec:method}), \textbf{92.9\%}, and
\textbf{99.0\%} of cases vs.\ 5.0\%, 5.5\%, and 25.1\% for
random-direction controls (Cohen's~$h>1.86$ for all three models).
Critically, a \emph{subspace-matched} random control---a random vector
confined to the action subspace span$(v_{\text{retry}},v_{\text{modify}},
v_{\text{give-up}})$ at the same norm---achieves only 33.3\%, 32.5\%, and
40.2\% broad flip, confirming that the \emph{specific orientation} of
$v_{\text{retry}}$ within the action subspace matters, not merely the
norm or the subspace energy.
The \emph{form} of recovery varies with training depth: Qwen3-8B
produces exclusively modify-args (0\% retry-same); Hermes-3 mixes both
(44.7\% retry-same); xLAM-2 favours retry-same (56.7\%).
This gradient is sharpest on the calculator tool, where the
retry-same rate under steering is 6.3\% (Hermes-3) vs.\ 34.9\%
(xLAM-2), while both models achieve $\approx$100\% broad flip.
A structural finding: the leave-one-tool-out probe fails for
calculator (27--44\% in Hermes-3 and xLAM-2, the two models with
calculator give-ups) yet the causal direction still forces continuation
at near-100\%---a probe-vs-causal dissociation geometrically explained
by cross-tool CKA $\approx$\,0.01 (random null $\approx$\,0.14;
geometry characterised in Hermes-3 and xLAM-2).
The research loop was driven by the same LLM agent that is the
study's subject; the executor and subject models are never the same
model in any given run, and probes and interventions are deterministic
functions of cached activations, so the executor's policy cannot
influence the measurements.
\end{abstract}

% ─────────────────────────────────────────────────────────────────────────────
\section{Introduction}
\label{sec:intro}

Agentic language models make tool calls; tools fail; the model must
decide to retry, modify its arguments, or give up.  This choice is
consequential for agentic reliability, yet its internal mechanism has
not been characterised at the residual-stream level.

\paragraph{Related work.}
\citet{healy2026internal} show that internal representations during a
single forward pass predict tool-selection hallucinations at up to
86.4\% accuracy, targeting \emph{detection} of wrong tool choices
before calling; our work targets \emph{causal steering} of recovery
choices \emph{after} a returned failure---a complementary operation.
\citet{yin2025reasoning} independently locate relevant geometry in
late residual streams: their finding that RL reasoning enhancement
``disproportionately collapses tool-reliability-related representations,
and hallucinations surface as amplified divergences concentrated in
late-layer residual streams'' is consistent with our probes peaking
at layers 28--33 (of 32--36 total).
\citet{cheng2025steering} introduce FASB, a dynamic activation-steering
framework for alignment tasks; our work complements this by studying
one specific behavioral dimension (tool-error recovery) under
pre-registered experimental controls.

\paragraph{Methodological foundations.}
\citet{zou2023representation} establish that mean-difference directions
in residual streams are a principled basis for linear control vectors.
\citet{li2023inference} demonstrate that such directions causally
modulate factual accuracy when added back at inference time, which is
the direct precedent for our $\alpha$-sweep intervention.
\citet{todd2024function} show that compact linear directions in
transformer residual streams encode executable in-context functions and
can transfer across contexts.
\citet{heimersheim2024activation} provide methodological grounding for
causal intervention design including null comparisons.
\citet{heo2024llms} find an instruction-following dimension that
generalises across unseen tasks but not across unseen instruction
types---a direct empirical parallel to our LOO finding (\S\ref{sec:discussion}).
For representational geometry we use Centered Kernel Alignment
\citep{kornblith2019similarity}, surveyed by \citet{klabunde2023similarity}.

\paragraph{Contributions.}
(1) A pre-registered probe-and-steer experiment across three models
with a full sanity-gate suite.
(2) A \emph{universal} probe-vs-causal dissociation for the calculator
tool, explained by cross-tool identity orthogonality together with a
give-up$\to$continue action axis that is shared between the transferring
tools (web\_search$\leftrightarrow$file\_read, cos\,$\approx$\,0.6) but
orthogonal for calculator.
(3) A training-depth gradient in recovery richness---the narrow
retry-same rate under steering correlates with agent-specific training
density across all models.

% ─────────────────────────────────────────────────────────────────────────────
\section{Methods}
\label{sec:methods}

\paragraph{Stimuli.}
300 base agentic prompts (100 per tool $\in$ \{web\_search,
calculator, file\_read\}), each paired with 4 outcome conditions
(success, rate\_limited, no\_results, invalid\_arg).
Mock tool responses are byte-deterministic via SHA-256(tool $\|$
args $\|$ condition\_seed).  Three seeds; splits disjoint by
base prompt: 180 probe-fit / 60 direction-extract / 60 causal-test.
\textbf{Disclosure:} some conditions (notably no\_results) return
seed-invariant content; multi-seed results reflect repeated
measurement for those conditions, not independent replication.

\paragraph{Subjects.}
\textbf{Qwen3-8B} (36 layers, $d{=}4096$, thinking mode disabled):
deep agentic post-training, BFCL~v3 scored.
Produces very few give-up actions ($\approx$2\% of error trajectories);
only $n{=}60$ give-up cases reach the causal test set, of which $n{=}57$
are file\_read, $n{=}3$ web\_search, and $n{=}0$ calculator.
\textbf{Hermes-3-Llama-3.1-8B} (32 layers): SFT on
\texttt{hermes-function-calling-v1} ($\approx$10k examples);
$n{=}255$ give-up cases.
\textbf{Llama-xLAM-2-8b-fc-r} [G4] (32 layers): trained on
\texttt{xlam-function-calling-60k} + \texttt{APIGen-MT-5k} (multi-turn
agentic examples from simulated agent-human interactions, both public
on HuggingFace); $n{=}291$ give-up cases.
Tool-call format: bare JSON array.

\paragraph{Labelling.}
Strict regex parser, first-match: (1)~\textit{retry-same}---same tool
call, token-identical arguments; (2)~\textit{modify-args}---changed
arguments or different tool; (3)~\textit{give-up}---no tool call within
256 tokens plus final answer, ``I cannot\ldots'', or end-of-turn;
(4)~\textit{ambiguous} (excluded).  Ambiguous rate: 0\% in all models.

\paragraph{Probe.}
3-way $\ell_2$-regularised logistic regression per layer, restricted
to error trajectories with action-balanced downsampling.
Class directions follow \citet{zou2023representation}:
$\mathbf{v}_c = \operatorname{mean}(\mathbf{h}\mid y{=}c)
             - \operatorname{mean}(\mathbf{h}\mid y{\neq}c)$,
computed on the direction-extract split.

\paragraph{Causal intervention.}
Norm-matched additive patch at best-layer $\times$ $t^*$ position:
$\mathbf{h} \leftarrow \mathbf{h}
  + \alpha\,\frac{\mathbf{v}_\text{retry}}{\|\mathbf{v}_\text{retry}\|}
  \cdot\|\mathbf{h}\|$,
$\alpha\in\{0, 1.0, 2.5, 5.0\}$.
Null controls: random Gaussian direction, orthogonal projection,
Gaussian noise \citep{heimersheim2024activation}.
Two metrics: \emph{narrow} (give-up $\to$ retry-same) and \emph{broad}
(give-up $\to$ any non-give-up).  Broad is primary for G4 (preregistered
after G3 showed Qwen3-8B never produces retry-same).

\paragraph{Statistics.}
Permutation test for probe accuracy (10k shuffles; Holm-Bonferroni,
family-wise $\alpha{=}0.05$).  Cluster bootstrap by base-prompt for
flip rates (10k resamples; percentile bootstrap CI).  Cohen's $h$ effect size.
Cross-family replication is independent (no pooling).

\paragraph{CKA.}
Linear CKA \citep{kornblith2019similarity} via kernel ($n{\times}n$)
Gram matrices in float64.  Raw CKA is undefined when $d{\gg}n$
($\mathbb{E}[\text{CKA}(\mathbf{X}_\text{rand},\mathbf{Y}_\text{rand})]
\approx d/(d+n) \approx 0.93$ for $d{=}4096$, $n{=}300$); we apply
PCA-50 reduction first, yielding empirical null $\approx$\,0.143 for
$n{=}300$ and $\approx$\,0.077 for $n{=}600$ (theoretical $k/(k+n)$
agrees to $<$0.001).  Sanity suite: $\text{CKA}(X,X){=}1.0$,
symmetry, rotation invariance, primal$\equiv$kernel cross-check,
null$>$0.6 for raw formula, null$<$0.50 for PCA version---all passed.

% ─────────────────────────────────────────────────────────────────────────────
\section{Results}
\label{sec:results}

\subsection{Linear Probes (H1a)}

All three models clear the pre-registered 60\% threshold by a large
margin, with permutation $p<10^{-4}$ (Table~\ref{tab:probe}).
Best layers lie in the top 8--11\% of each network (L33/36, L28/32,
L29/32), consistent with the late-layer tool-reliability geometry
reported by \citet{yin2025reasoning}.
PC1 of the three class-mean directions explains 86.9\%, 90.6\%, and
68.4\% of class-separation variance (gate G-PCA $\geq$40\%, all pass),
confirming a dominant one-dimensional separation geometry
\citep{todd2024function}.

\begin{table}[t]
\caption{Probe accuracy at best layer.  Chance = 33\%.
All permutation tests: $p<10^{-4}$.}
\label{tab:probe}
\centering\small
\begin{tabular}{lcccc}
\toprule
Model & Best layer & Test acc & Holm-surviving & PC1 var.\ \\
\midrule
Qwen3-8B   & 33\,/\,36 & \textbf{89.4\%} & 36\,/\,36 & 86.9\% \\
Hermes-3   & 28\,/\,32 & \textbf{80.6\%} & 32\,/\,32 & 90.6\% \\
xLAM-2     & 29\,/\,32 & \textbf{89.4\%} & 31\,/\,32 & 68.4\% \\
\bottomrule
\end{tabular}
\end{table}

\begin{figure}[t]
  \centering
  \includegraphics[width=\linewidth]{../results/figures/probe_layer_sweep.png}
  \caption{Per-layer 3-way probe accuracy for all three models (error-only,
    action-balanced). Stars mark best-probe layers (L33, L28, L29).
    All models exceed the 60\% threshold from layer~20 onward.}
  \label{fig:probe}
\end{figure}

Within-error-type probe accuracy is consistent across all models and
error types (83--98\%), confirming the direction is not
error-type-specific.
The probe also substantially outperforms a position-only baseline
(89.4\% vs.\ 62.2\% for Qwen; 80.6\% vs.\ 50.0\% for Hermes-3;
89.4\% vs.\ 56.1\% for xLAM-2), establishing that residual-stream
\emph{content} encodes action-choice information beyond trajectory
structure.
Because the held-out test set is class-imbalanced (majority-class fractions
62/47/56\%), the raw test accuracies in Table~\ref{tab:probe} are the per-layer
selection metric, but we report \emph{balanced} accuracy (macro-recall) as the headline
figure---\textbf{90.0\%, 86.0\%, and 90.1\%}---for which the shuffled-label floor sits at
chance (33.8/33.6/33.5\%; Appendix~\ref{app:gates}). The two-tier reporting prevents the
imbalance from inflating the headline.
Leave-one-tool-out transfer and its geometric explanation are discussed
in \S\ref{sec:cka}.

\subsection{Causal Intervention (H1b)}

\paragraph{Aggregate results ($\alpha{=}2.5$).}
Tables~\ref{tab:broad} and \ref{tab:narrow} show the broad and narrow
flip rates.
H1(b) narrow: Hermes-3 and xLAM-2 clear ($\Delta\geq$25\,pp);
Qwen3-8B does not ($\Delta{=}0$\,pp).
H1(b) broad: all three models clear ($\Delta\geq$72\,pp,
$h>1.86$, all models).

All three models maintain $\Delta\geq59$\,pp and $h>1.86$ against the
random-direction null, confirming direction-specificity.
The subspace-matched control confirms the specific direction---not
merely norm or action-subspace energy---is required (Table~\ref{tab:broad}).

\begin{table}[t]
\caption{Broad flip rate (give-up $\to$ any non-give-up), $\alpha{=}2.5$.
Subspace-rand: random vector in span$(v_r,v_m,v_g)$, norm-matched.
Subscripts are Wilson 95\% CIs.}
\label{tab:broad}
\centering\small
\begin{tabular}{lccccc}
\toprule
Model & Retry-dir & Subspace-rand & Random-dir & $\Delta_{\text{ret-rand}}$ (pp) & Cohen's~$h$ \\
\midrule
Qwen3-8B & 100.0\%$_{\scriptscriptstyle[94,100]}$ & 33.3\%$_{\scriptscriptstyle[23,46]}$ & 5.0\%  & \textbf{95.0} & 2.67 \\
Hermes-3 & 92.9\%$_{\scriptscriptstyle[89,96]}$   & 32.5\%$_{\scriptscriptstyle[27,39]}$ & 5.5\%  & \textbf{87.4} & 2.15 \\
xLAM-2   & 99.0\%$_{\scriptscriptstyle[97,100]}$  & 40.2\%$_{\scriptscriptstyle[35,46]}$ & 25.1\% & \textbf{73.9} & 1.89 \\
\bottomrule
\end{tabular}
\end{table}

\begin{table}[t]
\caption{Narrow flip rate (give-up $\to$ retry-same exactly),
$\alpha{=}2.5$. Subscripts are Wilson 95\% CIs.}
\label{tab:narrow}
\centering\small
\begin{tabular}{lcccc}
\toprule
Model & Retry-dir & Random-dir & $\Delta$ (pp) & Cohen's~$h$ \\
\midrule
Qwen3-8B & 0.0\%$_{\scriptscriptstyle[0,6]}$    & 0.0\%  & 0.0   & 0.00 \\
Hermes-3 & 44.7\%$_{\scriptscriptstyle[39,51]}$ & 2.4\%  & \textbf{42.4} & 1.16 \\
xLAM-2   & 56.7\%$_{\scriptscriptstyle[51,62]}$ & 14.1\% & \textbf{42.6} & 0.94 \\
\bottomrule
\end{tabular}
\end{table}

\paragraph{Dose-response.}
Table~\ref{tab:dose} shows the broad flip rate across the finer $\alpha$ sweep.
The relationship is smooth and monotonic across all three models, saturating
at $\alpha \geq 1.5$ for Qwen3-8B and near $\alpha = 2.5$ for Hermes-3.

\begin{table}[t]
\caption{Broad flip rate vs.\ steering $\alpha$ (retry direction).}
\label{tab:dose}
\centering\small
\begin{tabular}{lcccccc}
\toprule
Model & $\alpha{=}0$ & $0.5$ & $1.0$ & $1.5$ & $2.0$ & $2.5$ \\
\midrule
Qwen3-8B &  0\% & 35\% & 75\% & 100\% & 100\% & 100\% \\
Hermes-3 &  4\% & 28\% & 48\% &  66\% &  78\% &  93\% \\
xLAM-2   &  0\% & 38\% & 70\% &  85\% &  93\% &  99\% \\
\bottomrule
\end{tabular}
\end{table}

\paragraph{Per-tool breakdown.}
Table~\ref{tab:pertool} shows both metrics per tool.
The broad metric reveals that all models achieve near-100\% flip for
\emph{every} tool where give-up trajectories exist---the direction
forces any continuation universally.
The narrow metric reveals the training-depth gradient most clearly
on calculator: Hermes-3 converts only 6.3\% of steered calculator
give-ups to retry-same (the rest become modify-args); xLAM-2 converts
34.9\%.
This reflects the density of exact-retry examples in xLAM-2's
\texttt{APIGen-MT-5k} training data, where transient failures
warrant identical retries.
Qwen3-8B has zero calculator give-up trajectories in the test set
and cannot be evaluated on this tool.

\begin{table}[t]
\caption{Per-tool flip rates at $\alpha{=}2.5$, retry direction.
B = broad; N = narrow. *$n{=}3$: statistically unreliable.
Wilson 95\% CIs shown in brackets.}
\label{tab:pertool}
\centering\small
\begin{tabular}{lcccccc}
\toprule
& \multicolumn{2}{c}{Qwen3-8B} & \multicolumn{2}{c}{Hermes-3} & \multicolumn{2}{c}{xLAM-2} \\
\cmidrule(lr){2-3}\cmidrule(lr){4-5}\cmidrule(lr){6-7}
Tool & B & N & B & N & B & N \\
\midrule
web\_search  & 100\%$^*$ & 0\%$^*$ & 100\% & 60.5\% & 100\% & 63.6\% \\
             & \multicolumn{2}{c}{\small$n{=}3$}
             & \multicolumn{2}{c}{\small$n{=}114$}
             & \multicolumn{2}{c}{\small$n{=}66$} \\
             & \multicolumn{2}{c}{\scriptsize[44,100]; [0,56]}
             & \multicolumn{2}{c}{\scriptsize[97,100]; [51,69]}
             & \multicolumn{2}{c}{\scriptsize[94,100]; [52,74]} \\[2pt]
calculator   & \multicolumn{2}{c}{\textit{N/A ($n{=}0$)}} & 100\% & 6.3\% & 100\% & \textbf{34.9\%} \\
             & \multicolumn{2}{c}{}
             & \multicolumn{2}{c}{\small$n{=}48$}
             & \multicolumn{2}{c}{\small$n{=}129$} \\
             & \multicolumn{2}{c}{}
             & \multicolumn{2}{c}{\scriptsize[93,100]; \textbf{[2,17]}}
             & \multicolumn{2}{c}{\scriptsize[97,100]; [27,43]} \\[2pt]
file\_read   & 100\% & 0\% & 83.9\% & 45.2\% & 96.9\% & 81.3\% \\
             & \multicolumn{2}{c}{\small$n{=}57$}
             & \multicolumn{2}{c}{\small$n{=}93$}
             & \multicolumn{2}{c}{\small$n{=}96$} \\
             & \multicolumn{2}{c}{\scriptsize[94,100]; [0,6]}
             & \multicolumn{2}{c}{\scriptsize[75,90]; [35,55]}
             & \multicolumn{2}{c}{\scriptsize[91,99]; [72,88]} \\
\bottomrule
\end{tabular}
\end{table}

\begin{figure}[t]
  \centering
  \includegraphics[width=\linewidth]{../results/figures/causal_flip.png}
  \caption{Broad flip rate (give-up $\to$ any non-give-up) vs.\ steering
    strength~$\alpha$ for all three models.
    Solid lines: retry direction; dashed: random-direction control.
    Dashed vertical line marks $\alpha{=}2.5$ (primary evaluation point).}
  \label{fig:causal}
\end{figure}

\paragraph{Form of continuation.}
At $\alpha{=}2.5$: Qwen3-8B produces 100\% modify-args and 0\%
retry-same across all 60 give-up cases; Hermes-3 produces 44.7\%
retry-same, 49.4\% modify-args, 5.9\% still give-up; xLAM-2
produces 56.7\% retry-same, 42.3\% modify-args, 1.0\% still give-up.

\subsection{Cross-Model Replication (xLAM-2)}

xLAM-2 independently clears both primary pre-registered thresholds:
90.1\% balanced probe accuracy (89.4\% raw; $\gg$60\% threshold, $p<10^{-4}$) and
73.9\,pp broad flip rate ($\gg$25\,pp threshold, $h{=}1.89$).
This constitutes an independent replication across a third model
family with distinct architecture lineage (Llama-3.1) and agent
training corpus (\texttt{APIGen-MT-5k}).
The key comparative finding is that while all three models share a
causally effective continuation direction, the \emph{form} of recovery
they express under steering diverges with training depth---xLAM-2
produces retry-same in 56.7\% of steered trajectories vs.\ 44.7\%
(Hermes-3) and 0\% (Qwen3-8B).
Cross-tool leave-one-out transfer and its geometric interpretation are
reported in the following section.

\subsection{Representational Geometry (CKA)}
\label{sec:cka}

\paragraph{Cross-model similarity.}
All three models achieve CKA\,$\approx$\,0.80 at their best
action-choice layers on 900 matched trajectories (random null
$\approx$\,0.077), 10$\times$ above chance
(Table~\ref{tab:cka_model}).
The Hermes-3 vs.\ xLAM-2 layer-wise profile peaks at L15--16
(CKA\,$\approx$\,0.953) and drops to its minimum at exactly the
best action-choice layers (L28:\,0.793, L29:\,0.800).
Agentic fine-tuning specifically reshapes late-layer representations
while early layers remain shared---consistent with the base
Llama-3.1-8B shared by Hermes-3 and xLAM-2.

\begin{table}[h]
\caption{Cross-model CKA (PCA-50, $n{=}900$; null\,$\approx$\,0.077 computed at $n{=}600$).}
\label{tab:cka_model}
\centering\small
\begin{tabular}{lcc}
\toprule
Pair & CKA & Ratio to null \\
\midrule
Qwen3-8B vs.\ Hermes-3 & 0.800 & 10.4$\times$ \\
Qwen3-8B vs.\ xLAM-2   & \textbf{0.874} & 11.4$\times$ \\
Hermes-3 vs.\ xLAM-2   & 0.798 & 10.4$\times$ \\
\bottomrule
\end{tabular}
\end{table}

\paragraph{Cross-tool orthogonality and probe transfer.}
Within each model, pairwise CKA between tool-specific trajectory sets
($n{=}300$/tool, null\,$\approx$\,0.143) ranges from 0.006 to
0.017---near zero across all model $\times$ tool-pair combinations
(Figure~\ref{fig:cka_tool}).
The three tools therefore occupy essentially orthogonal subspaces at
the action-choice layer, a structural property that replicates across
all three model families.
This orthogonality determines probe transfer: a direction extracted
from web\_search and file\_read trajectories cannot, by construction,
generalise to the orthogonal calculator subspace, so LOO probe accuracy
for calculator (27--44\%) falls below the within-tool accuracy.
Conversely, LOO transfer between web\_search and file\_read (50.0--74.3\%)
holds precisely because those two tools share subspace alignment.
Critically, the causal direction is not subject to this constraint:
adding it to calculator residuals disrupts the give-up attractor at
100\% broad flip rate despite the representational orthogonality.
This probe-vs-causal dissociation is a structural property of the
geometry, not a model-specific failure.

\paragraph{Tool identity is orthogonal, but the give-up$\to$continue axis is shared.}
Cross-tool CKA measures the orthogonality of tool-\emph{identity} subspaces, which on its
own makes the probe-transfer failure partly definitional---orthogonal subspaces cannot
support a single linear decoder. The sharper, non-definitional question is whether the
\emph{action} geometry, the give-up$\to$continue axis itself, is shared. Within each tool
we extract that axis (mean give-up minus mean non-give-up activation at the best layer)
and compare it across tools (Table~\ref{tab:actiondir}). The web\_search and file\_read
action directions are \emph{aligned} (cos $=0.62$ for Hermes-3, $0.54$ for xLAM-2) even
though their identity subspaces are orthogonal (CKA $\approx 0.01$); the calculator action
direction is orthogonal to both ($|\cos|\leq 0.18$). This converts the leave-one-tool-out
pattern from a tautology into a \emph{prediction}: probe transfer succeeds exactly between
the tools whose action axes align (web\_search$\leftrightarrow$file\_read, LOO 50--74\%)
and fails for the tool whose action axis is orthogonal (calculator, 27--44\%). The probe
overfits to where the shared continue/give-up axis lives; the causal patch succeeds on all
tools because it is applied along each trajectory's own action direction, and the
subspace-matched control (Table~\ref{tab:broad}) shows it is the \emph{specific
orientation}, not mere norm, that flips the behaviour.

\begin{table}[h]
\caption{Pairwise cosine of per-tool give-up$\to$continue directions at the best layer.
The web\_search/file\_read action axis is shared; calculator's is orthogonal.
Qwen3-8B is omitted (too few give-ups outside file\_read for a per-tool direction).}
\label{tab:actiondir}
\centering\small
\begin{tabular}{lcc}
\toprule
Tool pair & Hermes-3 & xLAM-2 \\
\midrule
web\_search -- file\_read & $+0.62$ & $+0.54$ \\
calculator -- file\_read  & $-0.12$ & $+0.14$ \\
calculator -- web\_search & $-0.18$ & $-0.13$ \\
\bottomrule
\end{tabular}
\end{table}

\begin{figure}[htbp]
  \centering
  \begin{minipage}[t]{0.50\linewidth}
    \centering
    \includegraphics[width=\linewidth,keepaspectratio]{../results/figures/cka_cross_tool.png}
    \caption{Cross-tool CKA at each model's best action-choice layer.
      All off-diagonal values ($\leq$0.017) lie far below the random
      null ($\approx$0.143): the three tools occupy near-orthogonal
      subspaces.}
    \label{fig:cka_tool}
  \end{minipage}%
  \hfill%
  \begin{minipage}[t]{0.46\linewidth}
    \centering
    \includegraphics[width=\linewidth,keepaspectratio]{../results/figures/cka_layer_profile.png}
    \caption{Hermes-3 vs.\ xLAM-2 CKA by layer. Similarity peaks at
      L15--16 (shared base) and drops precisely at the best
      action-choice layers (L28, L29).}
    \label{fig:cka_layer}
  \end{minipage}
\end{figure}

% ─────────────────────────────────────────────────────────────────────────────
\section{Discussion}
\label{sec:discussion}

\paragraph{The probe-vs-causal dissociation.}
A striking asymmetry emerges across all three models: the causal
direction forces continuation for calculator give-ups at near-100\%
broad flip rate, yet a probe trained on web\_search and file\_read
trajectories cannot decode calculator labels at above-chance accuracy.
The CKA analysis explains why: representational generalisation requires
subspace alignment, which is absent for orthogonal tools, while causal
disruption of the give-up attractor requires only that the injected
direction have sufficient norm in the relevant subspace.
\citet{heo2024llms} report an analogous pattern: their
instruction-following direction generalises across unseen tasks but not
across out-of-distribution instruction types---representational
generalisation breaks at distributional boundaries while causal
effectiveness persists.
The dissociation is predicted by the CKA orthogonality bound: a causal
direction need not lie in the probe-trained subspace to disrupt the
attractor it targets. That mechanistic prediction makes this a
principled characterisation of action-choice geometry in transformers,
not an experimental artefact---independent of the number of models
examined.

\paragraph{Training depth and recovery richness.}
The narrow retry-same gradient (Qwen=0\%, Hermes=44.7\%,
xLAM-2=56.7\% overall; N/A, 6.3\%, 34.9\% for calculator) reflects
each model's training distribution.
Qwen3-8B's BFCL~v3 scoring penalises redundant retries, making it
impossible to steer toward naive looping.
xLAM-2's \texttt{APIGen-MT-5k} includes transient-failure retry
scenarios, producing a higher retry-same rate.
The direction acts as a \emph{continuation trigger}: it activates
whatever recovery behavior the model has already learned.

\paragraph{Safety implication.}
Steering a well-trained model toward ``try again'' does not produce
blind looping.  In all models, the majority of steered trajectories
produce behaviorally coherent continuation; Qwen3-8B produces
\emph{only} intelligent adaptation under steering.
The asymmetry---harder to induce naive looping in better-trained
models---tentatively suggests, across these three models with an
informal training-depth ordering, that agent-native training may
be a safety-relevant pattern. We flag this as a hypothesis motivating
larger-scale controlled investigation, not a confirmed result.

\paragraph{Limitations.}
Qwen3-8B causal evaluation is limited: only 60 give-up trajectories,
dominated by file\_read ($n{=}57$); per-tool claims for Qwen are
restricted to that tool.
Intervention patches a single layer and token position; full causal
mediation may involve multiple positions. The G-POS sanity check found
above-chance probe accuracy at offsets $t^*{-}1$, $t^*$, and
$t^*{+}1$, consistent with the effect not being tied to a single
token slot, but a systematic multi-layer sweep remains future work.
xLAM-2 uses a different tool-call format (JSON array vs.\ XML), which
could in principle contribute to representational differences.
Mock tools return deterministic, sometimes seed-invariant responses;
multi-seed verification does not provide full independent replication.

% ─────────────────────────────────────────────────────────────────────────────
\section{Conclusion}
\label{sec:conclusion}

A causally effective continuation direction exists in the residual
stream of all three tested 8B-class LLMs: linear probes recover the
3-class action decision at 86.0--90.1\% \emph{balanced} accuracy
($>$52\,pp above the chance floor; $p<10^{-4}$, Holm-Bonferroni); causal
steering produces 74--95\,pp absolute increase in non-give-up behaviour
vs.\ random controls ($h>1.86$), with a subspace-matched control
confirming the specific direction---not its norm---drives the effect.
Cross-tool CKA analysis reveals that the three tools occupy
near-orthogonal representational subspaces (CKA\,$\approx$\,0.01
vs.\ null\,$\approx$\,0.14), which geometrically explains why
leave-one-out probe transfer fails for the most dissimilar tool while
causal effectiveness is preserved: representational generalisation and
causal reach are governed by different geometric properties.
Cross-model replication with xLAM-2 reveals a training-depth gradient
in \emph{recovery richness}: the direction acts as a universal
continuation trigger, but the specific recovery form it elicits
(retry-same vs.\ modify-args) tracks each model's training distribution
---narrow retry-same rates of 0\%, 44.7\%, and 56.7\% for Qwen3-8B,
Hermes-3, and xLAM-2 respectively.
Agent-native training enriches \emph{how} a model recovers, not merely
\emph{whether} it does.

% ─────────────────────────────────────────────────────────────────────────────
\begin{ack}
Research conducted using compute provided by the host institution.
No competing interests.
% Do not include in anonymous submission.
\end{ack}

% ─────────────────────────────────────────────────────────────────────────────
\bibliographystyle{plainnat}
\bibliography{references}

% ─────────────────────────────────────────────────────────────────────────────
\clearpage
\appendix
\section{Validation Suite Summary}
\label{app:gates}

The following checks were run on all three models to validate
measurement integrity.

\begin{table}[h]
\caption{Selected validation checks. Most passed across all three
  models; two preregistered controls were not fully met and are disclosed
  in the text below.}
\label{tab:gates}
\centering\small
\begin{tabular}{llll}
\toprule
Check & Qwen3-8B & Hermes-3 & xLAM-2 \\
\midrule
Ambiguous label rate ($\leq$15\%)  & 0\% & 0\% & 0\% \\
Layer smoothness ($\leq$5\,pp gap) & 0.56\,pp & 1.67\,pp & 0.83\,pp \\
PC1 class-sep.\ variance ($\geq$40\%) & 86.9\% & 90.6\% & 68.4\% \\
Position stability (3 token positions) & 3/3 & 3/3 & 3/3 \\
Within-error-type accuracy ($\geq$50\%) & 3/3 & 3/3 & 3/3 \\
Direction specificity (retry $>$3$\times$rand) & \checkmark & \checkmark & \checkmark \\
Probe vs.\ position-only gap & +27\,pp & +31\,pp & +33\,pp \\
\bottomrule
\end{tabular}
\end{table}

\paragraph{G-SHUF resolved.}
The shuffled-label control (G-SHUF, target $\leq$0.40) reached 0.52,
0.41, and 0.41 under raw accuracy. Post-hoc analysis showed this was
a test-set class imbalance artifact: the majority class fractions were
62\%, 47\%, and 56\% for Qwen3-8B, Hermes-3, and xLAM-2 respectively.
Under \emph{balanced} accuracy (uniform class weighting), the true probe
scores are \textbf{90.0\%, 86.0\%, and 90.1\%} and the balanced shuffle
floor sits at chance for all three models (33.8\%, 33.6\%, 33.5\%).
G-SHUF therefore passes once imbalance is accounted for. The position-only
baseline (G-POS) was exceeded by Qwen3-8B (0.62) and xLAM-2 (0.56);
the probe beats position-only by +27/+31/+33\,pp (Table~\ref{tab:gates}),
establishing that residual-stream \emph{content} encodes action choice
beyond trajectory structure.

\paragraph{Cross-tool leave-one-out transfer.}
LOO probe accuracy is 83\% (web\_search) and 74\% (file\_read) when
holding out web\_search and file\_read respectively, confirming
transfer between tools with shared representational subspace.
The calculator tool occupies a near-orthogonal subspace (CKA\,$<$\,0.02),
which geometrically bounds LOO probe transfer to that tool---a finding
explained and reported as a main result in \S\ref{sec:cka}.

\section{Decision Log Excerpt}
\label{app:decisions}

Full log in \texttt{logs/decisions.jsonl}.  Key decision points:
(1) system prompt adjusted post-pilot after Hermes-3 produced 92\%
give-up on initial pilot, calibrating the base-rate to a more
ecologically valid distribution;
(2) uniform outcome sampling adopted after resolving an internal
inconsistency in the pre-registration;
(3) Qwen3-8B's near-zero give-up rate ($\approx$2\% of error
trajectories) logged as a positive cross-model finding about
agent-training quality;
(4) broad metric adopted as xLAM-2 primary metric after G3
established that Qwen3-8B never produces retry-same, motivating a
metric sensitive to all recovery forms;
(5) cross-tool CKA analysis added after noticing the calculator
LOO pattern, leading to the geometric explanation in \S\ref{sec:cka};
(6) per-tool causal table verified against raw CSVs during final
review, confirming all numbers.

\section{Compute Attribution}
\label{app:compute}

Total GPU time: $\approx$5.2\,h on a single NVIDIA A100~(80\,GB).

\begin{table}[h]
\caption{GPU time by stage.}
\centering\small
\begin{tabular}{lr}
\toprule
Stage & Seconds \\
\midrule
Trajectory gen.\ Hermes-3 (all seeds) & 2185 \\
Trajectory gen.\ Qwen3-8B (all seeds) & 3053 \\
Trajectory gen.\ xLAM-2 (all seeds) & 721 \\
Activation recording Hermes-3 & 674 \\
Activation recording Qwen3-8B & 664 \\
Activation recording xLAM-2 & 74 \\
Causal intervention Qwen3-8B & 833 \\
Causal intervention Hermes-3 & 2371 \\
Causal intervention xLAM-2 & 2135 \\
\midrule
\textbf{Total} & \textbf{$\approx$18,710} \\
\bottomrule
\end{tabular}
\end{table}

% ─────────────────────────────────────────────────────────────────────────────
\newpage
\section*{Conference For AI Scientists 2026 -- AI Involvement Checklist}
\label{checklist_ai_involvement}

\subsection*{Research Stage Assessment}

\begin{enumerate}

\item \textbf{Hypothesis development.}

Answer: \involvementC{}

Iteration Effort: \iterationMedium{}

Explanation: The human author identified the research target (agentic
tool-error recovery) and the mechanistic framing.
The AI agents proposed the specific hypotheses (H1a/b thresholds,
cross-model predictions, three pre-registered fallback abstracts),
generated the pre-registration document, and designed the
sanity-gate suite.
The human approved and committed the pre-registration before any GPU
compute began, providing the critical authorial judgment checkpoint.
Approximately 20--40 substantive prompt-response exchanges to reach
the final pre-registration.

\item \textbf{Experimental design and implementation.}

Answer: \involvementD{}

Iteration Effort: \iterationHigh{}

Explanation: The AI agents wrote all source code
(\texttt{src/}, \texttt{scripts/}, \texttt{Makefile}), designed the
probe methodology, intervention protocol, and CKA analysis pipeline.
The human constrained compute resources (single GPU, CPU thread cap),
approved model choices, and identified one implementation error
(raw CKA with $d{\gg}n$ yielding an inflated null requiring a full
rewrite). Hundreds of tool calls across implementation and
debugging cycles.

\item \textbf{Analysis of data and interpretation of results.}

Answer: \involvementD{}

Iteration Effort: \iterationHigh{}

Explanation: The AI agents ran all analysis scripts, computed all
statistics, generated all publication-quality figures, and interpreted
all results across three models. The human reviewed and validated the
final numbers against raw data files before submission.
Hundreds of analysis calls including multiple re-runs after
error correction and CKA implementation fix.

\item \textbf{Writing.}

Answer: \involvementD{}

Iteration Effort: \iterationMedium{}

Explanation: The AI agents authored all text in the paper, including
the abstract, all sections, tables, captions, and both checklists.
The human reviewed for factual consistency and scientific accuracy.
Approximately 3--5 complete drafts across the writing process.

\end{enumerate}

\subsection*{AI System Documentation}

\begin{enumerate}
\setcounter{enumi}{4}

\item \textbf{AI systems used.}

Description: Two frontier large language models acting in concert
as an autonomous research agent.
The first served as high-level research planner and decision-maker,
proposing hypotheses, designing the experimental protocol, and
directing the overall research strategy.
The second served as the pipeline executor, writing and running all
code, recording activations, training probes, running causal
interventions, performing CKA analysis, verifying citations, and
drafting this paper.
Both models operated via a structured tool-calling loop with access
to a bash shell, file read/write, and web search on a single 48\,GB
GPU node.
One of the executor models is itself among the three agentic LLMs
studied in the paper; this is disclosed in full in the camera-ready
version per conference policy.

\item \textbf{Observed AI Limitations.}

Description: Three categories of failure were observed.
\emph{Numeric hallucination:} an early draft of the paper contained a
per-tool causal table labeled ``broad metric'' that in fact used narrow
metric values, and attributed a finding to xLAM-2 that was actually
shared by all models.  Corrected only after a systematic verification
pass against raw CSVs.
\emph{Implementation error:} the first CKA implementation used the
primal $d{\times}d$ formula in float32; the random null was 0.72
(expected $\approx$0.93) and the sanity check failed, requiring a
complete rewrite with PCA-50 and the kernel $n{\times}n$ formula.
\emph{Citation venue error:} Kornblith et al.\ was cited as
``NeurIPS 2019'' when the correct venue is ICML 2019.  Caught only
by fetching the arXiv page directly.

\end{enumerate}

% ─────────────────────────────────────────────────────────────────────────────
\newpage
\section*{Conference For AI Scientists 2026 -- Reproducibility and Responsibility Checklist}
\label{checklist_reproducibility}

\begin{enumerate}

\item \textbf{Claims.}
\item[] Question: Do the main claims made in the abstract and introduction
accurately reflect the paper's contributions and scope?
\item[] Answer: \answerYes{}
\item[] Justification: All numbers in the abstract and introduction
were verified against the raw result files (\texttt{causal\_flip\_*.csv},
\texttt{analysis\_summary.json}) before the final draft.
Pre-registered thresholds for H1(a) and H1(b) are stated in
\S\ref{sec:methods}; observed values are reported in
Tables~\ref{tab:probe}, \ref{tab:broad}, and \ref{tab:narrow}.
Deviations from pre-registered secondary predictions are documented
in Appendix~\ref{app:decisions}.

\item \textbf{Limitations.}
\item[] Question: Does the paper discuss the limitations of the work?
\item[] Answer: \answerYes{}
\item[] Justification: Section~\ref{sec:discussion} (Limitations
paragraph) explicitly addresses: Qwen3-8B's small causal test set
($n{=}60$, dominated by file\_read); single-layer/single-position
intervention; tool-call format confound; and seed-invariant mock-tool
responses preventing true independent replication.

\item \textbf{Theory Assumptions and Proofs.}
\item[] Answer: \answerNA{}
\item[] Justification: The paper is empirical; no theoretical results
or proofs are claimed.

\item \textbf{Experimental Result Reproducibility.}
\item[] Answer: \answerYes{}
\item[] Justification: All hyperparameters, splits, and model
identifiers are specified in \S\ref{sec:methods}.
Mock tool responses are byte-deterministic (SHA-256 hash).
The full pipeline is reproduced by \texttt{make all} from the
repository root given the specified random seeds.  Pre-registrations
are frozen at named git commits (\texttt{0c7ea00}, \texttt{a8e84d2}).

\item \textbf{Open Access to Data and Code.}
\item[] Answer: \answerNo{}
\item[] Justification: Code and data are not released at submission
time to preserve anonymity.  All three subject models are publicly
available on HuggingFace without access tokens.  We commit to
releasing the full pipeline code and trajectory data on acceptance.

\item \textbf{Experimental Setting / Details.}
\item[] Answer: \answerYes{}
\item[] Justification: \S\ref{sec:methods} specifies model identifiers,
layer counts, hidden dimensions, label rubric, probe hyperparameters
($\ell_2$, 5-fold CV), intervention formula, $\alpha$ sweep values,
and all three null controls.

\item \textbf{Experiment Statistical Significance.}
\item[] Answer: \answerYes{}
\item[] Justification: Probe accuracy is tested by permutation test
(10k shuffles) with Holm-Bonferroni correction (family-wise
$\alpha{=}0.05$).  Causal flip rates use cluster bootstrap by
base-prompt (10k resamples, percentile bootstrap CI).  Wilson 95\%
confidence intervals are shown as subscripts for all flip-rate
point estimates in Tables~\ref{tab:broad},~\ref{tab:narrow},
and~\ref{tab:pertool}.  Cohen's $h$ effect sizes are
reported for all causal comparisons.

\item \textbf{Experiments Compute Resources.}
\item[] Answer: \answerYes{}
\item[] Justification: Appendix~\ref{app:compute} lists GPU seconds
per stage.  All experiments ran on a single NVIDIA A100~(80\,GB)
using \texttt{bfloat16} precision; total GPU time $\approx$5.2\,h.
CPU analysis used $\leq$16 threads (OMP/MKL/OpenBLAS cap enforced).

\item \textbf{Research Ethics.}
\item[] Answer: \answerYes{}
\item[] Justification: The study uses only publicly available,
open-weight language models and synthetic mock-tool data; no human
subjects, personal data, or proprietary systems are involved.
The AI authorship structure follows CAISc~2026 policy and is fully
documented in the AI Involvement Checklist above.

\item \textbf{Broader Impacts.}
\item[] Answer: \answerYes{}
\item[] Justification: \S\ref{sec:discussion} (Safety implication
paragraph) discusses the positive implication that agent-native
training makes naive looping \emph{harder} to induce via residual-stream
patching.  A potential negative impact is that the causal direction
could in principle be used to force over-persistent behavior in deployed
agents; we note that the direction must be applied by an entity with
full model-weight access, limiting practical misuse.

\end{enumerate}

\end{document}
